Til prioritering af features i spillet, anvendes MoSCoW-metoden, se \autoref{sec:moscowmetode}. Prioriteringerne kan ses i \autoref{fig:prioteringer}. Prioriteringerne er valgt, da det vurderes, at de tilsammen kan udgøre en hovedsløjfe, der er interessant og resulterer en samhørig oplevelse. Særligt afgrænses spillet også til spillertypen \textit{Achiever} jf. Bartles taksonomi, se \autoref{sec:bartletax}.
\begin{table}[H]
  \centering
\begin{tabular}{l p{10cm}}
\toprule
\textbf{Must have} &
\vspace*{-1em}
\begin{itemize}
\item Alle relevante oplægkrav*
\item PlayerAvatar
\item Pixel-Art
\item Skyde- og bevægelsessystem
\end{itemize} \\ \midrule

\textbf{Should have} &
\vspace*{-1em}
\begin{itemize}
\item Havstrøm-simulering
\item Flere antagonister
\item Controller-keyboard-support
\end{itemize} \\ \midrule

\textbf{Could have} &
\vspace*{-1em}
\begin{itemize}
\item Local Co-Op
\item Enorm diversitet af skraldsorter
\item Gysertematik
\item Meta-progression
\end{itemize} \\ \midrule

\textbf{Won't have} &
\vspace*{-1em}
\begin{itemize}
\item Multiplayer / Networking
\end{itemize} \\ \bottomrule
\end{tabular}\caption{I figuren ses prioteringerne fordelt baseret på vigtighedsgrad jf. \autoref{sec:moscowmetode}. *Alle relevante eksamenskrav kan ses i \autoref{sec:kravoplæg}}\label{fig:prioteringer}.
\end{table}
\subsubsection{Relevante oplægkrav}\label{sec:kravoplæg}
De relevante krav fra oplægget, der er obligatoriske er her opstillet og oversat. Jf. oplægget, se \autoref{sec:projektoplæg}.
\begin{enumerate}
    \item Brugeren skal undervises i korrekt håndtering af plastikaffald.
    \item I spillet anvendes økologisk simulering.
    \item Spillet skal indeholde en fjendtlig NPC med patrolscript.
    \item Adskillige assets skal være lavet fra bunden og udførlig begrundet sammenhæng den valgte målgruppe
    \begin{itemize}
        \item En målgruppe skal vælges (Under alle omstændigheder). 
    \end{itemize}
    \item Spillet skal itereres ad fem gange med baggrund i test af spil.
    \item Det endelige produkt skal bestå af henholdsvis en æstetisk og funktionel prototype. (\textbf{note} en prototype kan være begge)
    \item Produktet skal udgives interaktivt. (\textbf{note} kompilerede binærer)
    \item En filmet gennemspilning (playthrough) af spillet skal vedhæftes projektet.
\end{enumerate}
